%% Full elementary categorical proofs used by the lambda and internal-language sections.

\section{Limits, Yoneda, adjunctions, and exponentials}

\begin{annexitem}{annex:p10.cat.1}{thm:p10-limits-unique}{Uniqueness and the product equations}
\begin{proof}
Let \((L,\ell_j)\) and \((L',\ell'_j)\) be limits of \(D\).
Their universal properties give unique arrows
\(a:L\to L'\), \(b:L'\to L\) with
\(\ell'_j a=\ell_j\) and \(\ell_j b=\ell'_j\).
For every \(j\), \(\ell_j b a=\ell_j\).
Both \(b a\) and \(1_L\) therefore mediate the cone
\((L,\ell_j)\) to itself, so uniqueness gives \(b a=1_L\).
The same argument gives \(a b=1_{L'}\).
Any arrow preserving the displayed legs is the unique mediating
arrow already used, which proves the required uniqueness.
With an empty diagram the leg equations are vacuous, and the same
argument uses terminal uniqueness.

For colimits \((K,k_j)\) and \((K',k'_j)\), their universal
properties give \(a:K\to K'\), \(b:K'\to K\) with
\(a k_j=k'_j\) and \(b k'_j=k_j\).
The equations \(b a k_j=k_j\) and \(a b k'_j=k'_j\) give the two
identity composites by colimit uniqueness.  The mediating arrows
are unique with these leg equations, including the empty case.

If \(e:E\to A\) is an equalizer of \(f,g:A\to B\), and \(e r=e s\)
for \(r,s:X\to E\), then their common composite \(h:X\to A\)
equalizes \(f,g\).  Both \(r\) and \(s\) are factorizations of
\(h\) through the equalizer, hence are equal.
For a coequalizer \(q:B\to Q\), if \(r q=s q\), then the common
map out of \(B\) coequalizes the given parallel arrows.  The
coequalizer's unique factorization gives \(r=s\).

For product pairing, the projections of
\(\langle f,g\rangle a\) are \(f a,g a\), so the product property
identifies it with \(\langle f a,g a\rangle\).
The projections of \(1_A\times1_B\) are those of the identity
of \(A\times B\), giving the identity law.
The projections of \((f'\times g')(f\times g)\) are
\(f'f\pi_1,g'g\pi_2\), which are the projections of
\((f'f)\times(g'g)\).  Uniqueness proves the composition law.
\end{proof}
\end{annexitem}

\begin{annexitem}{annex:p10.cat.2}{thm:p10-limits-products-equalizers}{Constructing a diagram's universal cone and cocone}
\begin{proof}
Let \(D:J\to\mathcal C\) be a small diagram.  Form the two products
\[
 P=\prod_{j\in\operatorname{Ob}J}D j,\qquad
 Q=\prod_{u:j\to k\ \mathrm{in}\ J}D k.
\]
The second indexing set consists of all arrows of \(J\), including
its identities.  Write \(\pi_j:P\to D j\) and
\(\rho_u:Q\to D k\) for their projections.
There are unique \(s,t:P\to Q\) with
\[
                  \rho_u s=D(u)\pi_j,\qquad \rho_u t=\pi_k
                  \quad(u:j\to k).
\]
Take an equalizer \(e:L\to P\) of \(s,t\), and put
\(\ell_j=\pi_j e\).
Then
\(D(u)\ell_j=\rho_u s e=\rho_u t e=\ell_k\), so these are cone
legs.

Given a cone \(a_j:X\to D j\), product universality gives
\(a:X\to P\) with \(\pi_j a=a_j\).
For each \(u:j\to k\),
\(\rho_u s a=D(u)a_j=a_k=\rho_u t a\), hence \(s a=t a\).
The equalizer gives a unique \(b:X\to L\) with \(e b=a\).
This arrow satisfies \(\ell_j b=a_j\).
Conversely, any arrow satisfying these equations has
\(\pi_j e b=a_j\) for every \(j\), so \(e b=a\) by the product
property; equalizer uniqueness determines \(b\).
This proves the limit universal property.
If \(J\) is empty, \(P,Q\) are terminal, \(s=t\), and the
construction gives a terminal object as required.

For colimits, form
\[
 R=\coprod_{u:j\to k\ \mathrm{in}\ J}D j,\qquad
 S=\coprod_{j\in\operatorname{Ob}J}D j.
\]
Write \(\rho_u:D j\to R\) and \(\iota_j:D j\to S\) for
the injections.  Define \(s,t:R\to S\) by
\[
                       s\rho_u=\iota_k D(u),\qquad
                       t\rho_u=\iota_j.
\]
Take a coequalizer \(q:S\to K\) of \(s,t\), and set \(k_j=q\iota_j\).
Then \(k_k D(u)=q s\rho_u=q t\rho_u=k_j\), so these form a
cocone.
For any cocone \(a_j:D j\to X\), let \(a:S\to X\) be its
coproduct extension.  On each summand of \(R\),
\[
               a s\rho_u=a_k D(u)=a_j=a t\rho_u,
\]
and therefore \(a s=a t\).
The coequalizer gives a unique \(b:K\to X\) with \(b q=a\),
which has \(b k_j=a_j\).
If another map has the same composites with all \(k_j\), its
composite with \(q\) agrees with \(a\) on every coproduct summand,
so it equals \(a\), and coequalizer uniqueness gives the same \(b\).
For empty \(J\), \(R,S\) are initial and the construction yields
an initial object.

An empty product and binary products construct every finite
product recursively.  Repeated use of the binary universal
property gives the universal property for the entire finite
family.  Hence a terminal object, binary products, and
equalizers supply the first construction whenever \(J\) has
finitely many objects and arrows.  Initial objects and binary
coproducts similarly supply all finite coproducts and the
second construction.

In \(\mathsf{Set}\), the displayed compatible-tuples subset
has projections to the \(D j\)'s.  A cone of functions \(a_j:X\to D j\)
sends \(x\) to the tuple \((a_j(x))_j\), which is compatible by
the cone equations.  This is the only possible mediating
function, since every coordinate is prescribed.
For the colimit, the stated least equivalence relation exists
as the intersection of all equivalence relations containing
the generating pairs on the disjoint union.
The maps \(x\mapsto[(j,x)]\) form a cocone.
A cocone \(a_j:D j\to X\) induces the function
\([(j,x)]\mapsto a_j(x)\): the kernel equivalence relation of
the function on the disjoint union contains every generating
pair, so this formula is well-defined on the quotient.
It is unique because each quotient element has a representative
from some summand; no choice of representatives is required to
define the quotient function.
The product over the empty set is the singleton containing the
empty tuple, and the empty disjoint union and its quotient are
empty.  This proves the last assertions.
\end{proof}
\end{annexitem}

\begin{annexitem}{annex:p10.cat.3}{thm:p10-yoneda}{Evaluation and representable functors}
\begin{proof}
For \(x\in F(A)\), define
\(\alpha^x_X:\mathcal C(X,A)\to F(X)\) by
\(\alpha^x_X(f)=F(f)(x)\).
For \(u:Y\to X\), contravariant functoriality gives
\[
 F(u)\bigl(\alpha^x_X(f)\bigr)
       =F(u)F(f)(x)=F(fu)(x)=\alpha^x_Y(fu).
\]
Thus \(\alpha^x\) is natural.
At \(X=A,f=1_A\), functoriality gives
\(\alpha^x_A(1_A)=x\).

Conversely, if \(\alpha:\mathcal C(-,A)\Rightarrow F\) is
natural, apply its naturality square to the arrow \(f:X\to A\)
and the element \(1_A\in\mathcal C(A,A)\).
The result is
\[
                         \alpha_X(f)=F(f)(\alpha_A(1_A)).
\]
It determines every component of \(\alpha\), so the two
constructions are inverse.  They give a bijection between the
transformation family in the ambient universe and the small set
\(F(A)\); the parameter \(x\) determines all components, even when
the entire family of those components is large.

For naturality in the functor, let \(\theta:F\Rightarrow H\).
Evaluation of \(\theta\alpha\) at \(1_A\) is
\(\theta_A(\alpha_A(1_A))\), exactly the map induced by
\(\theta_A:F(A)\to H(A)\).
For naturality in the representing object, let \(h:A\to B\)
and \(\alpha:\mathcal C(-,B)\Rightarrow F\).
Precomposing \(\alpha\) with the transformation
\(\mathcal C(-,h):\mathcal C(-,A)\Rightarrow\mathcal C(-,B)\)
gives at \(1_A\) the element
\[
                    \alpha_A(h)=F(h)(\alpha_B(1_B)).
\]
This is the required naturality with the contravariant map
\(F(h):F(B)\to F(A)\).

Take \(F=\mathcal C(-,B)\).
The bijection identifies a transformation with an arrow
\(h:A\to B\), and its component at \(X\) is postcomposition
\(f\mapsto h f\).
Identity arrows induce identity transformations, and composing
two such transformations corresponds to composing their arrows.
This proves fullness and faithfulness.
If the transformation induced by \(h\) has an inverse, fullness
represents that inverse by an arrow \(k:B\to A\).
Faithfulness applied to the two identity composites gives
\(k h=1_A\) and \(h k=1_B\), proving the last assertion.
\end{proof}
\end{annexitem}

\begin{annexitem}{annex:p10.cat.4}{thm:p10-adjunction-data}{The hom-set bijections and the triangle identities}
\begin{proof}
Start with the natural bijections \(\Phi\), and define
\(\eta,\varepsilon\) as in the statement.
Naturality in the second variable, applied to \(1_{F A}\)
and \(f:F A\to B\), gives
\[
                         \Phi_{A,B}(f)=G(f)\eta_A.             \tag{1}
\]
Naturality in the first variable, applied to
\(\varepsilon_B:FGB\to B\) and \(g:A\to G B\), gives
\[
 \Phi_{A,B}(\varepsilon_B F(g))
       =\Phi_{G B,B}(\varepsilon_B)g=g.
\]
Consequently the inverse bijection is
\[
                         \Phi^{-1}_{A,B}(g)=\varepsilon_B F(g).
                                                                  \tag{2}
\]

For \(u:A\to A'\), naturality in the second variable at the
identity \(1_{F A}\) gives
\(\Phi_{A,F A'}(F(u))=GF(u)\eta_A\).
Naturality in the first variable at \(1_{F A'}\) gives the
same left side as \(\eta_{A'}u\).
Hence \(GF(u)\eta_A=\eta_{A'}u\), proving naturality of \(\eta\).
The inverse bijections are natural as well: solving the
naturality equation for \(\Phi\) through its inverse gives
\[
 \Phi^{-1}_{A',B'}(G(v)g u)
       =v\,\Phi^{-1}_{A,B}(g)\,F(u).
                                                                  \tag{3}
\]
For \(v:B\to B'\), use (3) first in the second variable at
\(1_{G B}\), and then in the first variable at \(1_{G B'}\).
Both evaluate \(\Phi^{-1}_{G B,B'}(G(v))\), giving
\[
                       v\varepsilon_B=\varepsilon_{B'}FG(v).
                                                                  \tag{4}
\]
This proves naturality of \(\varepsilon\).

Apply (2) to \(\eta_A=\Phi(1_{F A})\).
It yields \(\varepsilon_{F A}F(\eta_A)=1_{F A}\).
Apply (1) to \(\varepsilon_B=\Phi^{-1}(1_{G B})\).
It yields \(G(\varepsilon_B)\eta_{G B}=1_{G B}\).
These are the two triangle identities.

Conversely suppose natural \(\eta,\varepsilon\) with those
identities are given.  Define maps \(\Phi,\Psi\) by (1) and (2).
For \(f:F A\to B\), counit naturality and the first triangle give
\[
 \begin{aligned}
 \Psi(\Phi(f))
   &=\varepsilon_B\,FG(f)\,F(\eta_A)\\
   &=f\,\varepsilon_{F A}\,F(\eta_A)=f.
 \end{aligned}
\]
For \(g:A\to G B\), unit naturality and the second triangle give
\[
 \begin{aligned}
 \Phi(\Psi(g))
   &=G(\varepsilon_B)\,GF(g)\,\eta_A\\
   &=G(\varepsilon_B)\,\eta_{G B}\,g=g.
 \end{aligned}
\]
Thus these maps are inverse.
For \(u:A'\to A\) and \(v:B\to B'\), naturality of \(\eta\)
and functoriality give
\[
 \begin{aligned}
 \Phi_{A',B'}(v f F(u))
   &=G(v)G(f)GF(u)\eta_{A'}\\
   &=G(v)G(f)\eta_A u
    =G(v)\Phi_{A,B}(f)u.
 \end{aligned}
\]
Hence they define an adjunction.
Their values at the identity arrows are the given
\(\eta_A,\varepsilon_B\).
Equations (1) and (2) show that starting from any adjunction
bijections recovers those original bijections.
\end{proof}
\end{annexitem}

\begin{annexitem}{annex:p10.cat.5}{thm:p10-adjoint-preservation}{Transporting cones and cocones through an adjunction}
\begin{proof}
Let \(D:J\to\mathcal C\) have colimit \((L,\ell_j)\), and let
\(b_j:F(Dj)\to B\) be any cocone on \(F D\).
Its transposes
\[
                         a_j=\Phi_{D j,B}(b_j):D j\to G B
\]
form a cocone on \(D\).
Indeed, for \(u:j\to k\), the equations
\(b_k F(Du)=b_j\) and naturality in the first variable give
\(a_k D(u)=a_j\).
There is a unique \(a:L\to G B\) with \(a\ell_j=a_j\).
Let \(b=\Phi^{-1}_{L,B}(a):F L\to B\).
Naturality of the inverse bijections gives
\[
                  b F(\ell_j)=\Phi^{-1}_{D j,B}(a\ell_j)=b_j.
\]
If another arrow out of \(F L\) has these composites, its
transpose has composites \(a_j\) with all \(\ell_j\), hence
equals \(a\), and the bijection makes the original arrow \(b\).
This proves that \((F L,F\ell_j)\) is colimiting.

For the other assertion, let \(D:J\to\mathcal D\) have limit
\((L,\ell_j)\), and take a cone \(a_j:A\to G(Dj)\).
Its inverse transposes
\[
                         b_j=\Phi^{-1}_{A,D j}(a_j):F A\to D j
\]
form a cone on \(D\): naturality in the second variable
translates \(G(Du)a_j=a_k\) into \(D(u)b_j=b_k\).
Let \(b:F A\to L\) be their unique mediating arrow, and set
\(a=\Phi_{A,L}(b):A\to G L\).
Then
\[
                       G(\ell_j)a=\Phi_{A,D j}(\ell_j b)=a_j.
\]
Any other mediating arrow has an inverse transpose with the
same limiting cone legs, so equals \(a\).
Thus \((G L,G\ell_j)\) is limiting.

The proof also covers empty \(J\), where all leg conditions
are vacuous.  The listed constructions are the respective
diagram shapes in Definitions~\ref{def:p10-limit}
and~\ref{def:p10-products-equalizers}.
Their uniqueness is the compatible isomorphism of
Theorem~\ref{thm:p10-limits-unique}; no literal equality of
different choices was needed.
\end{proof}
\end{annexitem}

\begin{annexitem}{annex:p10.cat.6}{thm:p10-exponential-adjunction}{Transposition, function sets, and Heyting order}
\begin{proof}
For \(f:X\times A\to B\) and \(s:Y\to X\), the product
composition law and the defining evaluation equation give
\[
 \operatorname{ev}_{A,B}\bigl((\Lambda(f)s)\times1_A\bigr)
       =f(s\times1_A).
\]
Transpose uniqueness therefore gives
\(\Lambda(f(s\times1_A))=\Lambda(f)s\).

The product identity and composition laws make
\(X\mapsto X\times A,\ s\mapsto s\times1_A\) a functor.
For \(v:B\to B'\), define
\[
 v^A=\Lambda(v\operatorname{ev}_{A,B}):B^A\to(B')^A.
\]
Its defining equation is
\[
                 \operatorname{ev}_{A,B'}(v^A\times1_A)
                        =v\operatorname{ev}_{A,B}.          \tag{1}
\]
For \(v=1_B\), the identity of \(B^A\) satisfies this equation,
so transpose uniqueness gives \((1_B)^A=1_{B^A}\).
For \(w:B'\to B''\), composing (1) twice shows that
\(w^A v^A\) has untranspose \(wv\operatorname{ev}_{A,B}\);
uniqueness gives \((wv)^A=w^A v^A\).
Thus \(B\mapsto B^A\) is a functor.
For \(f:X\times A\to B\), equation (1) also gives
\[
                  \operatorname{ev}_{A,B'}(v^A\Lambda(f)\times1_A)
                        =v f.
\]
Hence \(\Lambda(v f)=v^A\Lambda(f)\).
Together with parameter naturality already proved, this is
the two-variable naturality of the exponential hom-set
bijections, establishing the adjunction.

In \(\mathsf{Set}\), a singleton admits exactly one function
from any set.  Cartesian products have the unique map
\(x\mapsto(f(x),g(x))\) specified by two component functions.
Let \(B^A\) be the set of functions from \(A\) to \(B\), and
put \(\operatorname{ev}(h,a)=h(a)\).
For \(f:X\times A\to B\), define
\(\Lambda(f)(x)\) to be the function \(a\mapsto f(x,a)\).
Evaluation recovers \(f\).
If \(k:X\to B^A\) has this evaluation, then for every \(x,a\),
\(k(x)(a)=f(x,a)\), so equality of functions gives
\(k=\Lambda(f)\).
When \(A\) is empty, \(B^A\) contains the unique empty function,
and both sides of the hom-set bijection are singletons.
The argument also includes empty \(B\) or \(X\); no existence
of a function from a nonempty set to the empty set is assumed.

For a Heyting algebra \(H\), its order defines the category
of the preceding body definition.
The greatest element \(1\) is terminal.
The meet inequalities and its greatest-lower-bound property
give the binary product \(a\wedge b\).
Residuation gives
\[
                        x\wedge a\leq b
                        \quad\Longleftrightarrow\quad
                        x\leq(a\Rightarrow b).
\]
Each corresponding hom-set is either empty or a singleton,
and this equivalence makes their cardinalities and their
unique possible arrows agree.
Evaluation is the inequality
\((a\Rightarrow b)\wedge a\leq b\), supplied by residuation
at \(x=a\Rightarrow b\).  Every arrow \(x\wedge a\to b\)
has the unique corresponding arrow \(x\to(a\Rightarrow b)\).
This is the exponential universal property, including
the one-element Heyting algebra.
\end{proof}
\end{annexitem}
